\section{Shared calibration and quantitative proof supplements}
\label{sec:v87-supplements}
The following statements supply details for the shared-apparatus and
meromorphic regimes above. They do not replace their hypotheses by
those of the projective spectrum experiment.

\subsection{Finite obstruction levels and continuity}
\begin{lemma}[Finite-level exclusion]\label{lem:v87-finite-level}
Let $Q$ be a nonempty topological space, let $Q_*\subset Q$ be dense,
and let $x:Q\to\R$ be continuous. If $x(Q_*)$ is contained in a finite
set $S$, then $x(Q)\subset S$. If $x$ is nowhere locally constant,
the preimage of each finite set is closed with empty interior, and
its complement is open and dense. If two continuous real-valued
functions agree on that complement, they agree everywhere.
\end{lemma}
\begin{proof}
Since $S$ is closed, $x^{-1}(S)$ is closed and contains the dense set
$Q_*$, and therefore equals $Q$. For any $s\in S$, an interior point
of $x^{-1}(s)$ would give an open set on which $x$ is constant.
Thus every $x^{-1}(s)$ is closed with empty interior. A finite union
of such sets again has empty interior: intersecting their open dense
complements gives an open dense set, by a finite induction.
Finally the equality set of two continuous real-valued functions is
closed, so agreement on a dense set implies agreement everywhere.
\end{proof}
In Theorem~\ref{thm:v86-shared}, different clock constants would force
the true action on the regular set into at most one direct obstruction
level and two swapped obstruction levels. The lemma makes the
contradiction explicit. After the constants agree, it extends equality
of the actions off the finitely many exceptional levels to all of $Q$.
The choice between the two labels may vary from point to point; no
continuation of eigenvectors, no local constancy of that choice, and
no assumption excluding arbitrarily fine label oscillation is used.

\subsection{A quantitative certificate for four shared anchors}
The four-site calibration procedure has a finite, computable separation
certificate. It displays separately the conditioning of an admissible
label choice and the distance to a false calibration.

Fix four sites as in Proposition~\ref{prop:v86-calibration}, with four
distinct actions and noncoincident two-clock profiles. Choose one
representative $(r_{i1},r_{i2})$ of the recovered odds at each site.
For $\sigma\in\{0,1\}^4$ set
\[
 s_{ij}^{\sigma}=\begin{cases}r_{ij},&\sigma_i=0,\\
                              r_{ij}^{-1},&\sigma_i=1,
             \end{cases}\qquad
 A_\sigma=(s_{i1}^{\sigma},-s_{i2}^{\sigma})_{i=1}^4,
 \quad b=(T_1-T_2)\one.
\]
Let $H$ be a fixed compact rectangle contained in $(0,\infty)^2$
whose interior contains $h_*=(1/k_1,1/k_2)$. Define
\begin{align}
 \mathcal A&=\{\sigma:A_\sigma h_*=b\},\nonumber\\
 \lambda&=\min_{\sigma\in\mathcal A}\sigma_{\min}(A_\sigma),
 \qquad d=\min_{\sigma\notin\mathcal A}\min_{h\in H}
                                    \|A_\sigma h-b\|_2.
                                               \label{eq:v87-anchor-certificate}
\end{align}
An empty outer minimum in $d$ is interpreted as $+\infty$.

\begin{proposition}[Certified four-anchor calibration]
\label{prop:v87-anchor-certificate}
The constants $\lambda,d$ in \eqref{eq:v87-anchor-certificate} are
positive. Suppose all recovered odds and their inverses are bounded
by a fixed number $M$, and suppose perturbed representatives can be
matched to them, independently at each site, with entrywise error at
most $\varepsilon\le(2M)^{-1}$. There is $C=C(M,H)$ such that a minimizer
\[
 (\widehat\sigma,\widehat h)\in
       \mathop{\rm argmin}_{\sigma\in\{0,1\}^4,\ h\in H}
                                  \|\widehat A_\sigma h-b\|_2
\]
satisfies
\begin{equation}\label{eq:v87-anchor-bound}
 C\varepsilon<d/3\quad\Longrightarrow\quad
 \widehat\sigma\in\mathcal A,\qquad
                     \|\widehat h-h_*\|_2\le C\varepsilon/\lambda.
\end{equation}
The clock constants $k_j=1/h_j$ have the same local error order.
The procedure does not require the true label choices as input.
\end{proposition}
\begin{proof}
The obstruction argument for four distinct action values proves that
any positive solution of any $A_\sigma h=b$ must equal $h_*$:
a different calibration could account for at most three distinct
levels. This argument uses positivity, not a prior bound on a
competing action range. The true label pattern belongs to $\mathcal A$.
If an active matrix had a nonzero null vector, an arbitrarily small
movement of $h_*$ in that direction would stay in the interior of
$H$ and give another positive solution, a contradiction. Hence
$\lambda>0$. Inactive matrices have no solution on $H$; compactness
and the finite number of patterns give $d>0$.

For sufficiently small $\varepsilon$, inversion is Lipschitz on
the fixed odds interval, so all sixteen matrix perturbations have
operator norm at most $C\varepsilon$. At an active choice and $h_*$
the perturbed residual is at most $C\varepsilon$; the minimum has
no larger residual. Since $H$ is bounded, the unperturbed residual
of a minimizing pair is therefore at most $C\varepsilon$, after
increasing $C$. If $C\varepsilon<d/3$, this pair is active. Now
\[
 \lambda\|\widehat h-h_*\|_2
 \le\|A_{\widehat\sigma}(\widehat h-h_*)\|_2
 =\|A_{\widehat\sigma}\widehat h-b\|_2
 \le C\varepsilon.
\]
The reciprocals of the coordinates of $H$ have bounded derivatives,
which proves the final assertion.
\end{proof}
The certificate involves only sixteen two-variable constrained
least-squares problems once the noiseless odds are fixed. Bounds for
nearby exact odds follow by the same finite minimization. For the
correct direct labelling, the determinant of two rows of $A$ has
absolute value
\[
        k_1k_2|T_2-T_1|\,|x_i-x_l|.
\]
Thus repeated anchor levels and small residuals of false swapped
calibrations are distinct sources of poor conditioning. To start
from noisy \emph{matrices}, one must also control the odds-extraction
error at the four regular anchors; on the compact regular classes
of Proposition~\ref{prop:v86-calibration} it is bounded by a constant
times the matrix error. This proposition is a finite-anchor certificate,
not a claim of uniform Lipschitz recovery through a two-clock profile
collision or a minimax theorem for an unrestricted spatial field.

\subsection{The binary determinant cancellation, coefficient by coefficient}
In the proof of Theorem~\ref{thm:v86-product-modulus}, write
$e=h'-h$, $C_b=U_{:b}V_{:b}^{\mathsf T}$ and
$\mathfrak d=\det U\det V$. With $r=a/(1+a)$,
\[
 K(T)=s(T)C_1+t(T)C_2,
 \qquad s(T)=r(T-x),\quad t(T)=(1-r)(T-\tau).
\]
Put $H=(T_1-h)(T_2-h)$. The affine interpolant of the true matrices
is $Q=q_1C_1+q_2C_2$, where
\begin{align*}
 q_1(T)&=r\left[1+(h-x)(T_1+T_2-h-T)/H\right],\\
 q_2(T)&=(1-r)\left[1+(h-\tau)(T_1+T_2-h-T)/H\right].
\end{align*}
All four affine scalar polynomials have uniformly bounded coefficients.
Since $\det(vC_1+wC_2)=vw\mathfrak d$, the exact difference is
\begin{equation}\label{eq:v87-explicit-det-cancellation}
 \det(K-eQ)-\det K
       =\mathfrak d\{-e(sq_2+tq_1)+e^2q_1q_2\}.
\end{equation}
Every coefficient is at most $C\rho(|e|+e^2)$ in absolute value.
The pole estimate $|e|\le C\delta/\gamma$, the bounded pole interval,
and $\rho\le C\gamma$ make this at most $C\delta$.

For the second perturbation, let an affine matrix be $M_0+TM_1$.
Its determinant coefficients are
\[
 \det M_0,\qquad
 (M_0)_{11}(M_1)_{22}+(M_1)_{11}(M_0)_{22}
 -(M_0)_{12}(M_1)_{21}-(M_1)_{12}(M_0)_{21},\qquad
 \det M_1.
\]
On $\|M_0\|_F+\|M_1\|_F\le R$, the difference of each coefficient
under a coefficient perturbation of size $\nu\le1$ is at most
$C(R+1)\nu$, by subtracting the displayed products. Here
$K'-(K-eQ)$ is the interpolant of
$(T_j-h')(P'_j-P_j)$, $j=1,2$, whose coefficient norm is at most
$C\delta$. No inverse channel factor is introduced by this second
step. This proves the required $C\delta$ determinant-coefficient
bound before division by its leading coefficient of order $\rho$.

\subsection{Finite-cell radii}
For binary observations, twelve entries across the three clocks are
controlled. If every entry error is at most $t$, each Frobenius error
is at most $\sqrt4\,t=2t$. Thus the radius
$\eta=\sqrt{(2/n)\log(24/\zeta)}$ satisfies
\[
 \Pr\{\max_j\|\widehat P_j-P_j\|_F>\eta\}
             \le24\exp(-n\eta^2/2)=\zeta.
\]
This supplies explicitly the norm conversion in
Theorem~\ref{thm:v86-adaptive}; independence between different cells
is not needed for the union bound.

\subsection{Periods and entire branched fibres}
In Lemma~\ref{lem:v86-normalized}, imaginary-period normalization
means vanishing real parts of periods on the \emph{punctured}
surface. A basis of the $2g$ handle cycles together with small residue
loops generates those periods, subject to the usual one relation
between puncture loops. A small positive loop around a pole of real
residue $r$ has period $2\pi i r$, already imaginary. The holomorphic
correction fixes the handle periods; there is no additional omitted
condition at a puncture. Multiplication by $i$ converts this convention
to real periods.

For Theorem~\ref{thm:v86-branched}, the finite set
$S_\pi=\pi(\operatorname{supp}D)\cap I_0$ is only a candidate set.
At a point $p$ over $s$, of ramification index $e_p$, a nonzero
secant charge produces a simple pole of $\dd\pi/(\pi-s)$ with
residue $e_p$, whereas the tangent $\dd\pi/(\pi-s)^2$ has pole
order $e_p+1$. Consequently the sharper support candidates are
\[
 S_{\rm sec}=\{s\in I_0: (\pi^{-1}(s))_{\rm red}\le D\},
 \qquad
 S_{\rm tan}=\{s\in I_0:\sum_{p\mapsto s}(e_p+1)p\le D\}.
\]
They are finite subsets of $S_\pi$. The charge test may be restricted
to $S_{\rm sec}$ and the tangent test to $S_{\rm tan}$: outside them
membership in $E$ is impossible. The proof is the disjointness of
fibres and the displayed local pole orders. These support conditions
are necessary, not substitutes for membership in the actual subspace
$E$. They leave the stated sufficient clock budgets unchanged.
